\begin{abstract}
Light carries data and signatures of the path and objects it encounters,
allowing the same optical infrastructure to support communication and sensing.
Optical integrated sensing and communication (O-ISAC) builds on this dual
capability across fiber, free space optical, visible light, and light fidelity
(LiFi) platforms. Photonic hardware extends the concept to millimeter wave and
terahertz systems, while hybrid designs connect physical domains. This breadth
motivates comparison of the design insights that extend
across platforms and those tied to a particular physical setting. A systematic
literature search covering 1 January 2020 through 22 June 2026 identified 227
eligible reports representing 206 unique studies. We organized this evidence by
physical platform and coupling mechanism using structured narrative synthesis
and a technical assessment developed for this review. The comparison framework
connects each result to its task, measurement plane, operating conditions, and
validation setting. Across the included corpus, bandwidth, optical power,
hardware constraints, and signal design jointly shape communication reliability
and sensing performance. These relationships depend on reported conditions and
extend beyond a single rate and resolution frontier. Laboratory studies span
diverse architectures. Broader field validation, shared benchmarks, and
reusable data and code emerge as priorities for stronger comparison and
reconstruction. The resulting taxonomy and framework support the design and
evaluation of O-ISAC systems for sixth generation (6G) networks.
\end{abstract}

\begin{IEEEkeywords}
6G, integrated sensing and communication, metric comparability, optical
communication, optical sensing.
\end{IEEEkeywords}
